\section{Problem Formulation and Notation}\label{sec:formulation}

\subsection{Notation}\label{sec:notation}
$\mathbf q=[q_0,q_1,q_2,q_3]^\top\in\mathbb{S}^3$ is a unit quaternion and $\mathbf R(\mathbf q)\in SO(3)$ the corresponding rotation. For a 3-vector $\mathbf p$, define the lifted monomial vector
\begin{equation}
\mathbf v_2(\mathbf q)=\big[q_0^2,q_1^2,q_2^2,q_3^2,\sqrt2 q_0q_1,\dots,\sqrt2 q_2q_3\big]^\top\in\mathbb{R}^{10},
\label{eq:v2}
\end{equation}
so that $\|\mathbf v_2(\mathbf q)\|_2^2=\|\mathbf q\|_2^4=1$ for $\mathbf q\in\mathbb{S}^3$, and $\mathbf e_{\rm tr}=[1,1,1,1,0,\dots,0]^\top$ satisfies $\mathbf e_{\rm tr}^\top\mathbf v_2(\mathbf q)=1$.

\textbf{Fact 3.1 (affine representability).} For every $\mathbf p\in\mathbb{R}^3$ and $a\in\{1,2,3\}$ there is $\mathbf m_a(\mathbf p)\in\mathbb{R}^{10}$ with $\mathbf e_a^\top\mathbf R(\mathbf q)\mathbf p=\mathbf m_a(\mathbf p)^\top\mathbf v_2(\mathbf q)$ for all $\mathbf q\in\mathbb{S}^3$ (verified to $4\times10^{-15}$); this is why all four constraint families are affine in $(\mathbf v_2,\mathbf t)$.

\subsection{Objective and constraints}\label{sec:objective}
\begin{equation}
J_\theta(\mathbf q,\mathbf t)=\mathbf v_2(\mathbf q)^\top\mathbf M_0(\theta)\mathbf v_2(\mathbf q)+2\mathbf t^\top\mathbf D_0(\theta)\mathbf v_2(\mathbf q)+\mathbf t^\top\mathbf G_0(\theta)\mathbf t .
\label{eq:objective}
\end{equation}
For PnP-type residuals $\mathbf r_i(\mathbf q,\mathbf t)=\mathbf P_i(\mathbf R(\mathbf q)\mathbf p_i+\mathbf t)$ with $\mathbf P_i=\mathbf I-\mathbf u_i\mathbf u_i^\top/\|\mathbf u_i\|^2$ the bearing projection, \eqref{eq:objective} holds with $\mathbf M_0=\sum_i\sum_{ab}P_{i,ab}\operatorname{sym}(\mathbf m_a(\mathbf p_i)\mathbf m_b(\mathbf p_i)^\top)$, $\mathbf D_0=\sum_i\mathbf P_i[\mathbf m_1;\mathbf m_2;\mathbf m_3]$ and $\mathbf G_0=\sum_i\mathbf P_i$. This is QPEP's problem class minus the elimination. All constraints are of the form $g_j(\mathbf q,\mathbf t)=\tilde{\mathbf a}_j^\top\mathbf v_2(\mathbf q)+\mathbf b_j^\top\mathbf t\le0$; Table~\ref{tab:families} lists the four families.

\begin{table}[t]
\centering
\caption{Constraint families. All four were verified to $4\times10^{-15}$. The tilt-cone row uses $\|\mathbf q\|^2=1$, i.e.\ it is $q_1^2+q_2^2\le\sin^2(\alpha/2)$ on the manifold.}
\label{tab:families}
\scriptsize
\setlength{\tabcolsep}{2.5pt}
\begin{tabular}{@{}p{1.9cm}p{3.55cm}p{2.55cm}@{}}
\toprule
Family & Expression & $(\tilde{\mathbf a}_j,\mathbf b_j)$ \\
\midrule
Positive depth & $z_{\min}-\mathbf e_3^\top(\mathbf R\mathbf p_i+\mathbf t)\le0$ & $(-\mathbf m_3(\mathbf p_i)+z_{\min}\mathbf e_{\rm tr},\ [0,0,-1]^\top)$ \\
Planar non-pen. & $-\mathbf n^\top(\mathbf R\mathbf v_j+\mathbf t)-d\le0$ & $(-\sum_a n_a\mathbf m_a(\mathbf v_j)-d\,\mathbf e_{\rm tr},\ -\mathbf n)$ \\
Field of view ($\pm$) & $\pm(\mathbf R\mathbf p_i+\mathbf t)_x-\tau(\mathbf R\mathbf p_i+\mathbf t)_z\le0$ & $(\pm\mathbf m_1(\mathbf p_i)-\tau\mathbf m_3(\mathbf p_i),\ [\pm1,0,-\tau]^\top)$ \\
Tilt cone & $q_1^2+q_2^2-\sin^2(\tfrac\alpha2)\|\mathbf q\|^2\le0$ & $(\mathbf e_2+\mathbf e_3-\sin^2(\tfrac\alpha2)\mathbf e_{\rm tr},\ \mathbf 0)$ \\
\bottomrule
\end{tabular}
\end{table}

The two FoV inequalities imply $z>0$ only if $\tau<\infty$ and the constraint is active in the right direction, so we always add one explicit positive-depth constraint per observed point; this is cheap ($p$ grows by the number of points) and removes any separate chirality argument.

\subsection{Problem (C-QPEP)}\label{sec:cqpep}
\begin{equation}
\min_{\mathbf q\in\mathbb{S}^3,\ \mathbf t\in\mathbb{R}^3}\ J_\theta(\mathbf q,\mathbf t)\quad\text{s.t.}\quad g_j(\mathbf q,\mathbf t;\theta)\le0,\quad j=1,\dots,p. \tag{C-QPEP}
\end{equation}
\begin{assumption}\label{asm:main}
$\mathbf M_0,\mathbf D_0,\mathbf G_0$ are continuously differentiable in $\theta$, and $\mathbf G_0\succ0$ (or the active $\mathbf b_j$ span $\mathbb{R}^3$).
\end{assumption}

\textbf{Out of scope.} (i) Non-affine constraints such as $\|\mathbf R\mathbf p+\mathbf t-\mathbf c\|\ge d$ or inter-object clearance. (ii) Multiple optima exchanging as $\theta$ varies (the solution map jumps). (iii) Non-polynomial dependence on $\mathbf q$, e.g.\ a geodesic distance to a reference rotation.
